% !TeX spellcheck = en_US

\section{Conclusions and Future Work} \label{sec:concl}

This final section synthesizes the main scientific and methodological conclusions derived from the development, optimization, and validation of the \textit{AstroPT MM} architecture. It also outlines the prospects and future lines of research established by this thesis in the field of foundation models applied to astrophysics.

\subsection{Scientific and Methodological Conclusions} \label{sec:concl:sci}

The core achievement of this thesis is the successful design and optimization of \textit{AstroPT MM Autoregressive Opt}, a bimodal foundation model (yet natively prepared for scaling to additional modalities) of $\sim 100.75$ million parameters. This model integrates, for the first time under a self-supervised autoregressive framework, spatial photometric information from \textit{Euclid} and fine spectroscopic information from \textit{\gls{DESI}}. Pretraining was completed on a curated catalog of $137\,661$ unique galaxies, unifying optical and near-infrared images from \textit{Euclid} (DR1) with optical spectra from \textit{\gls{DESI}} (\textit{\gls{DESI}} DR1). Based on the systematic evaluation and quantitative and qualitative astrophysical analysis presented in \cref{sec:res}, the following fundamental conclusions are drawn:

\begin{enumerate}
	\item \textbf{Multimodal Synergy and Breaking Observational Degeneracies:} 
	Evaluation via \textit{downstream tasks} demonstrates that the joint latent representation space robustly outperforms individual unimodal baselines across the majority of the evaluated physical parameters. The joint model achieves a good coefficient of determination in estimating first-order physical parameters, such as total stellar mass ($\log M_*$, $R^2 = 0.9339$) and star formation rate (\textit{log \gls{SFR}}, $R^2 = 0.7487$), significantly exceeding the predictive capabilities of both the isolated visual and spectral branches. 
	
	Physically, this systematic increase in joint-space accuracy confirms that unsupervised bimodal latent fusion effectively breaks degeneracies intrinsic to isolated instrumental domains. By coupling high-resolution morphology with 1D spectra, the latent space successfully resolves reddening and age-metallicity degeneracies. This establishes that the multimodal foundation model does not merely concatenate data, but organically learns to intersect spatial structure with chemical composition to infer global galaxy properties.
	
	\item \textbf{Competitiveness of Continuous Autoregression:} 
	The continuous autoregressive pretraining of \textit{AstroPT MM Autoregressive Opt} provides a robust and competitive framework compared to traditional contrastive alignment schemes. In direct comparison with \textit{AstroCLIP} \autocite{ar:astroclip}, our model matches or improves upon its performance in the joint modality for the vast majority of physical and geometric tasks. 
	
	Crucially, in the pure visual channel, \textit{AstroPT MM} outperforms \textit{AstroCLIP} on key physical targets such as stellar mass ($R^2 = 0.8427$ versus $0.7526$) and \textit{\gls{SFR}} ($R^2 = 0.6398$ versus $0.6037$). This demonstrates that autoregressive learning with token mixing compels the encoder to model high-fidelity morphological features instead of relying on low-frequency global alignments.
	
	\item \textbf{Outperforming the Unimodal Supervised Limit:} 
	One of the critical findings of this project is that the joint-space embeddings of the \textit{AstroPT MM Autoregressive Opt} architecture systematically outperform task-specific unimodal supervised reference models (such as the convolutional \textit{ResNet18} network for images and the Conv1D+Transformer model for spectra). 
	
	Although the latter baseline models were trained with direct optimization on physical labels for a single task, the bimodal pretraining of \textit{AstroPT MM Autoregressive Opt} instills a density of physical knowledge and synergistic cross-modal interconnection that more than compensates for the lack of direct supervision. This performance gap is fundamentally explained by two distinct factors. First, for integrated global parameters (such as stellar mass and \gls{SFR}), the visual modality provides crucial chromatic structural information that complements the central spectrum, acting as an effective proxy for the mid-infrared \textit{\gls{WISE}} photometry used to compute the \gls{VAC} ground truth. Second, for pure spectroscopic parameters (such as emission lines), where the visual modality cannot provide relevant physical information, the outperformance is strictly driven by the scale and diversity of the pretraining data. Ultimately, this strongly validates the foundation model hypothesis in observational astrophysics: unsupervised representation learning over massive unlabeled data yields richer and highly generalized latent representations than traditional supervised paradigms restricted by human label availability.
	
	\item \textbf{Limitations of Hybrid Contrastive Alignment:} 
	The architectural variant \textit{Hybrid AstroPT MM}, which introduces InfoNCE contrastive losses at the midpoint of the transformer (layer 6) using a dual-phase attention mask, reported significantly lower performance in estimating physical variables ($R^2 = 0.8675$ for joint-modality stellar mass and $0.6447$ for \textit{\gls{SFR}}). 
	
	This experimental behavior shows that imposing a global vector similarity bottleneck at the midpoint of the transformer's flow severely restricts the non-linear expressiveness that the bimodal autoregressive model freely exploits from the early layers. Contrastive alignment degrades the detailed physical representation of individual physical flux patches, limiting its ability to model complex transitions on the multimodal latent manifold.
	
	\item \textbf{Self-Attention Dynamics and the $F(\text{H}\beta)$ Anomaly:} 
	The performance drop experienced by the nebular recombination line flux $\text{H}\beta$ in the joint space ($R^2 = 0.7374$ compared to the excellent $R^2 = 0.9594$ of the isolated spectroscopic branch) highlights an important limitation in the bimodal optimization dynamics of causal self-attention. This anomaly indicates that the integration of spatial features may inadvertently introduce noise or suppress specific fine-grained spectral features during joint pretraining, emphasizing the need for future work to investigate and balance the loss functions to preserve isolated spectral fidelity.
	\item \textbf{Absence of Modal Transfer:} As noted in \cref{sec:res:cross}, knowledge transfer between modalities during training is deficient. While some appreciable transfer exists when comparing the \textit{Autoregressive Opt} model against \textit{SSL Unimodal I} or \textit{Supervised I} (or their spectral counterparts), it is still far from achieving cross-modal transfer that matches performance in tasks where they are not the dominant modality. Improving this deficit will be fundamental, especially to obtain spectroscopic information from images acquired with next-generation surveys.
\end{enumerate}

In conclusion, this work consolidates continuous autoregressive pretraining of bimodal patches as a viable and highly efficient paradigm for constructing unified representation models of the local universe.

\subsection{Future Work and Scaling Perspectives} \label{sec:concl:future}

Future research lines and experimental extensions derived from this work are structured along the following technological and astrophysical axes:

\begin{itemize}
	\item \textbf{Value-Added Catalogs via Latent Transfer:} 
	The ultimate, albeit currently unresolved, observational goal of \textit{AstroPT MM} is the high-precision inference of physical parameters for galaxies lacking spectroscopic data. While the current architecture exhibits a deficit in cross-modal transfer, solving this limitation in future iterations would allow the visual embeddings to absorb and retain spectral knowledge during pretraining. This would theoretically make it possible to construct accurate estimators for photometric redshifts ($z_{\text{phot}}$), stellar masses, and \textit{\gls{SFR}} from pure optical/near-infrared images, unlocking the creation of massive synthetic catalogs without the need for dedicated spectroscopic follow-up.
	
	\item \textbf{Unsupervised Scientific Applications in Observational Astronomy:} 
	As documented in \cref{sec:ap:sim}, the richness of the bimodal space opens the door to direct unsupervised exploitation. The $k$-NN similarity search engine implemented on the joint latent space enables ultra-fast searches to identify and catalog rare astrophysical sources of high scientific relevance, such as \textit{dual \glspl{AGN}}.
	
	\item \textbf{Generative Models and Latent Diffusion in Embedding Spaces:} 
	One of the most innovative proposals, experimentally introduced in \cref{sec:ap:diff}, is the development of conditioned latent diffusion models within the latent space of \textit{AstroPT MM Autoregressive Opt}. By training a diffusive model to transition between visual morphology embeddings and the continuous flux spectrum, the direct synthesis of optical spectra conditioned on the visual morphological features of \textit{Euclid} will be enabled. This tool will act as a ``virtual spectrograph'', predicting emission lines and resolved metallicities from wide-field photometry.
	
	\item \textbf{Data Scaling and Modeling in the Chinchilla Limit:} 
	Since the current data volume places pretraining in a suboptimal token regime relative to the model parameters, future work contemplates incorporating the complete overlap with future massive surveys. Currently, the model is prepared to support new data modalities from ground-based surveys, such as those carried out by the \textit{\gls{HSC}} project or the \textit{\gls{SDSS}} survey. Incorporating wide-field survey data from the \textit{Vera C. Rubin Observatory} (\textit{\gls{LSST}}) and the future \textit{Nancy Grace Roman Space Telescope} will provide the volume of objects necessary to reach the optimal convergence zone and classical scaling laws. It is expected that a larger data volume and a richer pretraining set will enable the model to continue improving physical parameter predictions.
	
	\item \textbf{\gls{MoE} Architectures:} 
	To scale the bimodal foundation model beyond 100 million parameters without drastically increasing the computational cost on \textit{TeideHPC}, we propose restructuring the network using a state-of-the-art architecture such as \textit{\gls{MoE}}. In this scheme, specific transformer layers will act as experts specialized in either spatial morphology or fine-grained spectral kinematics, dynamically routing visual and spectral tokens to optimize the physical representation of the joint space.
	
		
	\item \textbf{Exploration of Discrete Tokenizers in Large Data Volume Regimes:} 
	Although the initial discretization experiment using the pretrained encoders of \textit{AION-1} proved unfeasible due to quantization noise and instrumental domain mismatch, this approach retains significant potential value for future work with larger-scale samples. 

	With the future incorporation of massive datasets from complete surveys (such as the \textit{Euclid} DR1 releases and the commissioning of the \textit{\gls{LSST}}), we plan to develop and pretrain native bimodal discrete tokenizers from scratch (using techniques such as \textit{\gls{FSQ}} or \textit{\gls{LFQ}}). A critical data volume of several million galaxies will allow vector quantizers to structure dense and robust codebooks, capable of retaining subtle surface brightness gradients and fine spectral indicators without suffering information loss from quantization collapse, effectively enabling the treatment of the observable universe as a unified discrete language.
	
	\item \textbf{Unified Pretraining Strategy on Unpaired Data:} 
	One of the greatest physical bottlenecks for scaling \textit{AstroPT MM} is the restriction of training exclusively on samples with direct overlap (the positional matching catalog of $137\,661$ objects). To bypass this data volume limit, we propose a scalable pretraining strategy based on severe modality masking. Under this semi-paired scheme, the model can ingest unpaired data by simply masking $100\%$ of the tokens belonging to the missing modality. The pretraining process would alternate between pure unimodal batches (e.g., passing a sequence of \textit{Euclid} patches followed by completely masked \textit{\gls{DESI}} patches) and fully paired multimodal batches using standard token mixing. This allows the architecture to seamlessly assimilate millions of isolated images or spectra without altering the core autoregressive mechanism, exploiting the entirety of the data available in each survey independently.
\end{itemize}